\documentclass[11pt]{article}
\usepackage[margin=1in]{geometry}
\usepackage{amsmath,amssymb,amsthm}
\usepackage{xurl}
\usepackage{hyperref}
\sloppy
\let\oldus\_ \renewcommand{\_}{\oldus\allowbreak}
\newtheorem{theorem}{Theorem}
\newtheorem{lemma}[theorem]{Lemma}
\theoremstyle{definition}
\newtheorem{definition}[theorem]{Definition}
\newcommand{\ind}{\operatorname{ind}}

\title{Conjecture 00000001673 is false:\\ the inducibility of $C_6$ is at most $2/7 < (\sqrt3-1)/2$}
\author{Jackmeson1}
\date{}

\begin{document}
\maketitle

\section*{The conjecture}
\emph{Definition: the inducibility of $C_6$ is its maximal asymptotic density as an induced subgraph
of $n$-vertex graphs. Conjecture: $\ind(C_6) = (\sqrt3-1)/2$; the extremal configuration is the
balanced blow-up of the triangular prism, uniquely; and uniqueness is verified by the stabilizer
structure of triangular-prism blow-ups.}

We refute the value clause (A) for every graph on $n \ge 7$ vertices, and hence for the limit,
the limsup and the liminf of the maximum induced density, and for the limsup along any sequence of
graphs. We also refute the extremal-configuration clause (B): every blow-up of the triangular prism
has no induced $C_6$ at all, while $\ind(C_6) \ge 5/324 > 0$. The clause about the ``stabilizer
structure'' is not a mathematical statement we formalize; it presupposes (B).

\section*{Definitions and reading}
\begin{definition}[Induced copy]
For a graph $G$ and a $6$-set $S \subseteq V(G)$, $S$ \emph{induces} $C_6$ if $S$ is the vertex set
of an induced embedding of the $6$-cycle, i.e.\ an injective $f\colon \mathbb{Z}/6 \to V(G)$ with
$f(i)f(j) \in E(G) \iff j-i \equiv \pm 1 \pmod 6$. Equivalently $G[S] \cong C_6$.
\end{definition}

\begin{definition}[Density and inducibility]
For $G$ on $n$ vertices, $d(G) = \#\{S : |S|=6,\ S \text{ induces } C_6\} / \binom{n}{6}$, and
$M(n) = \max_{|V(G)|=n} d(G)$. Following Pippenger and Golumbic, as quoted by Hatami, Hirst and Norine
(arXiv:1108.5699): ``Let [$i_H(n)$] denote the maximum induced density of [$H$] in any graph on
[$n$] vertices. ... The asymptotic behavior of [$i_H(n)$] is captured by the inducibility of [$H$]
defined as [$\lim_{n\to\infty} i_H(n)$]. It is straightforward to see that this limit always
exists.'' (Bracketed symbols were lost in text extraction.) The Lean development defines
$\ind(C_6) := \limsup_{n} M(n)$ and also treats $\liminf_n M(n)$, any limit of $M(n)$, and
$\limsup_n d(G_n)$ for an arbitrary sequence $(G_n)$.
\end{definition}

\begin{definition}[Prism and blow-ups]
The triangular prism is $K_3 \,\Box\, K_2$ (``the Cartesian product of a cycle graph with a single
edge'', Wikipedia, \emph{Prism graph}; ``Triangular prism graph -- 6 vertices, 9 edges''). A
blow-up of $H$ along a map $p\colon V \to V(H)$ has $uv \in E \iff p(u)p(v) \in E(H)$; it is balanced
when all fibres $p^{-1}(x)$ have equal size (Hatami--Hirst--Norine: ``the copies of two vertices are
adjacent if and only if the originals are. If every vertex is replaced with the same number of
copies, then the resulting graph is called a balanced blow-up.'').
\end{definition}

Conventions: graphs are simple and loopless; $\mathbb{N}$ includes $0$; densities are real numbers;
$\sqrt{3}$ is the positive root. No normalisation other than $\binom n6$ is used.

\section*{Proof of (A)}
\begin{lemma}[Facts about $C_6$]
$C_6$ is $2$-regular, has no twins (for $i \neq j$ some $k \notin\{i,j\}$ is adjacent to exactly one of
$i,j$), and is triangle-free.
\end{lemma}

\begin{lemma}[Pair lemma]
Let $b$ be a vertex set and $x \ne y$ in $b$ such that $b\setminus\{x\}$ and $b\setminus\{y\}$
both induce $C_6$. Then every $w \in b\setminus\{x,y\}$ satisfies $w \sim x \iff w \sim y$.
\end{lemma}
\begin{proof}
Let $S = b\setminus\{x,y\}$. In the copy $b\setminus\{y\} = S \cup\{x\}$ the vertex $w$ has
degree $2 = |N(w)\cap S| + [w\sim x]$; in $b\setminus\{x\} = S\cup\{y\}$ it has degree
$2 = |N(w)\cap S| + [w\sim y]$. Hence $[w\sim x] = [w \sim y]$.
\end{proof}

\begin{lemma}[Local bound]
For every vertex set $b$, at most two $x \in b$ have $b\setminus\{x\}$ inducing $C_6$.
\end{lemma}
\begin{proof}
Suppose $x,y,z$ are distinct and all three deletions induce $C_6$. By the pair lemma applied to $y,z$,
the vertices $y,z$ have the same neighbours in $b\setminus\{y,z\}$. But $y,z$ both lie in the induced
copy $b\setminus\{x\}$, and since $C_6$ has no twins there is $w \in b\setminus\{x\}$, $w\ne y,z$, adjacent
to exactly one of $y,z$. Contradiction.
\end{proof}

\begin{theorem}[Main bound]
For every $n \ge 7$ and every graph $G$ on $n$ vertices, $d(G) \le 2/7$.
\end{theorem}
\begin{proof}
Let $c$ be the number of $6$-sets inducing $C_6$. Double count pairs $(S,T)$ with $S$ such a
$6$-set, $|T| = 7$, $S \subset T$. Each $S$ lies in exactly $n-6$ sets $T = S\cup\{x\}$; each $T$ contains
at most $2$ such $S$ (each $S \subset T$ is $T\setminus\{x\}$; local bound). So
$c(n-6) \le 2\binom n7$. With $7\binom n7 = (n-6)\binom n6$ and $n-6>0$ this gives
$7c \le 2\binom n6$.
\end{proof}

Since $(11/7)^2 = 121/49 < 3$, $\sqrt3 > 11/7$, so $(\sqrt3-1)/2 > 2/7$. Therefore $M(n) \le 2/7$ for
$n\ge7$, hence $\limsup M$, $\liminf M$, any limit of $M$, and $\limsup d(G_n)$ along any sequence are
all $\le 2/7 < (\sqrt3-1)/2$. In particular (A) fails in every one of these readings.

\section*{Proof of (B)}
\begin{theorem}
No blow-up of the triangular prism (balanced or not) contains an induced $C_6$.
\end{theorem}
\begin{proof}
Let $f$ be an induced embedding of $C_6$ into the blow-up along $p$, and $g = p\circ f$. If
$g(i) = g(j)$ with $i\ne j$, then $f(i), f(j)$ have identical adjacency to every $f(k)$, so $i,j$ are
twins in $C_6$, impossible. So $g$ is injective into the $6$ prism vertices, hence bijective, and
$i\sim j \iff g(i)g(j) \in E(\text{prism})$. The prism triangle $(0,0),(1,0),(2,0)$ then pulls back to
a triangle in $C_6$, impossible.
\end{proof}

\begin{theorem}[Lower bound]
For $m\ge1$, $M(6m) \ge 5/324$; hence $\ind(C_6) = \limsup M \ge 5/324$.
\end{theorem}
\begin{proof}
In the blow-up of $C_6$ on $\{0,\dots,6m-1\}$ along $v\mapsto v \bmod 6$, each $\varphi\colon \{0..5\}\to\{0..m-1\}$
gives the induced copy $i\mapsto 6\varphi(i)+i$, and distinct $\varphi$ give distinct vertex sets. So
$c \ge m^6$, while $\binom{6m}{6} \le (6m)^6/6!$, giving $d \ge 720/6^6 = 5/324$.
\end{proof}
Hence every sequence of prism blow-ups has density identically $0$, which does not tend to
$\ind(C_6) \ge 5/324$, and at every $n = 6m$ the maximum density exceeds theirs by at least $5/324$.

\section*{Lean formalization}
File \texttt{lean/Conjecture1673/Basic.lean} (Lean 4, Mathlib \texttt{v4.33.1}, only \texttt{import Mathlib}).
Key definitions: \texttt{C1673.InducesC6} (exists an embedding \texttt{f : cycleGraph 6 -> G} of type \texttt{SimpleGraph.Embedding} with image $S$),
\texttt{c6Count}, \texttt{c6Density}, \texttt{maxC6Density} (\texttt{iSup} over all graphs on \texttt{Fin n}),
\texttt{indC6} (\texttt{limsup maxC6Density atTop}), \texttt{prism} (\texttt{K3 box K2} via \texttt{SimpleGraph.boxProd}),
blow-ups as \texttt{prism.comap p}.
Key theorems: \texttt{pair\_twins}, \texttt{erase\_count\_le\_two}, \texttt{count\_mul\_le},
\texttt{c6Density\_le}, \texttt{two\_sevenths\_lt}, \texttt{indC6\_le}, \texttt{liminf\_le}, \texttt{lim\_le},
\texttt{seq\_limsup\_le}, \texttt{prism\_blowup\_no\_C6}, \texttt{prism\_c6Density}, \texttt{indC6\_ge}, and the
main theorem \texttt{C1673.not\_conjecture}, which states the conjunction of: density $<(\sqrt3-1)/2$ for all
$n\ge7$; \texttt{indC6} $\ne(\sqrt3-1)/2$; $\liminf \ne (\sqrt3-1)/2$; $M(n) \not\to (\sqrt3-1)/2$; for every graph
sequence $\limsup d(G_n) \ne (\sqrt3-1)/2$; no sequence of prism blow-ups has density tending to \texttt{indC6};
prism blow-ups on $6m$ vertices have density at least $5/324$ below $M(6m)$.

Axiom audit: \texttt{\#print axioms C1673.not\_conjecture} reports
\texttt{[propext, Classical.choice, Quot.sound]}; no \texttt{sorry}, no \texttt{native\_decide}. The finite
facts about $C_6$ are checked by kernel \texttt{decide} on \texttt{Fin 6}.

\end{document}
